\section{Introduction}
\IEEEPARstart{A}{I} accelerators are increasingly communication-bound. While arithmetic throughput continues to grow through parallel processing elements, tensor cores, and specialized neural processing units, overall system efficiency is often limited by how quickly weights, activations, partial sums, embeddings, and key-value cache data can move through the memory hierarchy and between compute tiles. This trend shifts the design focus from isolated compute blocks toward the communication substrate that connects them \cite{dally2001route,chen2016eyeriss}.

In this context, the Network-on-Chip (NoC) has evolved from a scalable replacement for buses and crossbars into a broader system fabric for AI chips. The same communication backbone must now support conventional on-die traffic, one-to-many data reuse in deep neural networks, inter-die movement across chiplets, trusted execution, and, in some forward-looking systems, hybrid optical links. The result is that latency, throughput, energy per bit, and trustworthiness all depend on interconnect design choices rather than on compute architecture alone \cite{bjerregaard2006survey,weerasena2024security}.

This evolution is driven by real industrial trends. Over the past five years, AI accelerator design has shifted decisively toward heterogeneity and modularity. Recent commercial accelerators from NVIDIA, AMD, Intel, and Tesla increasingly adopt multi-die or chiplet organizations connected through advanced packaging and high-bandwidth die-to-die links. At the same time, the emergence of large language models has pushed model sizes beyond what a single reticle-sized die can accommodate, making chiplet-based scaling a practical necessity rather than a design option. On a longer horizon, silicon photonics is being explored as a way to break the energy-per-bit barrier of electrical links for global on-package communication. These developments mean that the NoC is no longer just a chip-level detail; it is a system-level architectural decision that cuts across compute, memory, packaging, security, and software.

This report follows the structure emphasized in our group presentation while also following the project guidance to focus primarily on trends from roughly the past three to five years and to extract design insight rather than merely classify papers. Section~\ref{sec:fundamentals} introduces NoC fundamentals. Section~\ref{sec:aiopt} discusses AI-specific optimizations such as multicast-aware routing and compression. Section~\ref{sec:chiplet} explains why heterogeneous integration and chiplets make interconnect design even more critical. Section~\ref{sec:security} analyzes NoC security and trustworthiness in heterogeneous AI systems. Section~\ref{sec:photonic} surveys emerging photonic NoC directions. Section~\ref{sec:crosscut} then compares these directions under a unified architectural lens. We conclude by arguing that NoC design should be treated as a first-class architectural problem for scalable AI computing.
